% Generated draft: requires mathematical and source review.
\documentclass[11pt]{article}
\usepackage[margin=25mm]{geometry}
\usepackage{fontspec}
\setmainfont{Latin Modern Roman}
\usepackage{amsmath,amssymb,mathtools,unicode-math}
\setmathfont{Latin Modern Math}
\usepackage{enumitem,needspace,xurl,hyperref}
\setlength{\parindent}{0pt}
\setlength{\parskip}{5pt}
\setlength{\emergencystretch}{3em}
\newcommand{\N}{\mathbb{N}}
\newcommand{\Z}{\mathbb{Z}}
\newcommand{\Q}{\mathbb{Q}}
\newcommand{\R}{\mathbb{R}}
\newcommand{\C}{\mathbb{C}}
\newcommand{\F}{\mathbb{F}}
\newcommand{\A}{\mathbb{A}}
\newcommand{\PP}{\mathbb{P}}
\newcommand{\E}{\mathbb{E}}
\newcommand{\eps}{\varepsilon}
\newcommand{\dd}{\,\mathrm{d}}
\begin{document}
\section*{UnsolvedMath ID: 30003313}
\textbf{Source problem code:} catalogue entry
\subsection{Definitions and mathematical statement}
Let \(\mathbf{No}\) denote the proper class of surreal numbers with its usual ordered-field structure. A derivation on a field \(K\) is an additive map \(D:K\to K\) satisfying \(D(ab)=aD(b)+bD(a)\). A differential field is a field equipped with a derivation. Two differential fields are elementarily equivalent in a specified common first-order language if they satisfy the same sentences in that language. The source does not specify the additional conditions intended by a \emph{surreal derivation}.

The supplied record contains a collection of six questions. Several essential definitions and hypotheses are absent, so the following preserves the questions and identifies the information required for precise formulations.
\begin{enumerate}
\item Does there exist a satisfactory composition operation on all of \(\mathbf{No}\)? The source does not specify the admissible pairs of surreal numbers, the axioms for composition, its required compatibility with existing operations, or the foundational interpretation of an operation on this proper class. In particular, ``on all of \(\mathbf{No}\)'' does not determine whether composition must be defined for every pair of surreal numbers.

\item Are the field of EL-series and the omega-series field elementarily equivalent as differential fields? Symbolically, the requested comparison has the form
\[
(K_{\mathrm{EL}},D_{\mathrm{EL}})\equiv(K_{\omega},D_{\omega}).
\]
Here these symbols name the structures mentioned in the source; their constructions, coefficient fields, chosen derivations, and precise comparison language are not supplied. This notation therefore does not resolve which structures are intended.

\item Does every \(\kappa\)-computable function on the Cauchy completion of \(\mathbf{No}_{\kappa}\) satisfy the intermediate value theorem? Write \(C_{\kappa}\) provisionally for the completion intended by the source. The usual intermediate value property for a function between linearly ordered domains is the condition
\[
\begin{gathered}
\forall a,b\in\operatorname{dom}(f)\quad\forall y\in\operatorname{codom}(f),\\
a<b\ \land\ \min\{f(a),f(b)\}<y<\max\{f(a),f(b)\}\\
\Longrightarrow\ \exists c\in\operatorname{dom}(f)\quad
 a<c<b\ \land\ f(c)=y.
\end{gathered}
\]
This displays the conventional property only; the source does not determine the domain and codomain needed to instantiate it. It also leaves unspecified the permitted values of \(\kappa\), the meaning of \(\mathbf{No}_{\kappa}\), the notion of Cauchy completion, and the definition and representation conventions for \(\kappa\)-computability. No continuity hypothesis is supplied.

\item Is there a Kaplansky differential embedding theorem for \(\mathbb{R}((G))_{\kappa}\)? Ordinarily, \(\mathbb{R}((G))\) denotes a Hahn field of formal sums with real coefficients and well-ordered support in an ordered abelian group \(G\). The source does not define the subscript \(\kappa\), specify hypotheses on \(G\) or \(\kappa\), choose a derivation, or state the source fields, target fields, and properties that the proposed embedding theorem should involve. Consequently, a quantified embedding assertion cannot be recovered from this record alone.

\item Determine a function, or characterize functions, \(f\) satisfying
\[
\lambda_{f(x)}=\kappa_x.
\]
The maps or families denoted by \(\lambda\) and \(\kappa\), the domain of \(x\), the domain and codomain of \(f\), and the scope of the identity are unspecified. The original wording asks for ``\(f\)'', whereas the rewritten wording asks which functions satisfy the identity. Neither existence, uniqueness, nor an intended class of solutions is specified. No identification of the notation \(\kappa_x\) with the parameter \(\kappa\) in the other questions is justified by the supplied record.

\item Do there exist distinct surreal derivations \(D_1,D_2\) on \(\mathbf{No}\) such that
\[
D_1\ne D_2,
\qquad
D_1\big|_{\mathbb{R}\langle L\rangle}
 =D_2\big|_{\mathbb{R}\langle L\rangle}?
\]
The source does not define \(L\), the substructure \(\mathbb{R}\langle L\rangle\), or the additional axioms for surreal derivations. This question occurs in the original statement but is omitted from the supplied cleaned and current statements; it is retained here pending review of the intended scope.
\end{enumerate}
\subsection{Short English statement}
This record groups questions about composition of surreal numbers, elementary equivalence of differential series fields, an intermediate value property for computable functions, differential embeddings of Hahn fields, an identity involving two unspecified families, and distinct surreal derivations agreeing on a specified substructure. Missing definitions and hypotheses prevent fully determined formulations. The sixth question is preserved from the original statement despite its omission from the rewritten statements. The supplied open label is a dated curation claim, not an independently verified conclusion.
\subsection{Sources}
\begin{itemize}
\item Local UnsolvedMath registry: supplied record with permanent ID \(30003313\), title \emph{Composition and Elementary Equivalence in Surreal Differential Fields}, exact source problem code \texttt{catalogue entry}, in the Oberwolfach Reports Open Problems collection. No individual proposer is supplied; the proposed year is recorded as 2016.
\item Supplied report citation: \emph{Mini-Workshop: Surreal Numbers, Surreal Analysis, Hahn Fields and Derivations} (2017), Oberwolfach Reports, pp.~3313--3372, DOI \texttt{10.4171/owr/2016/60}. \url{https://doi.org/10.4171/owr/2016/60}. This is also the supplied literature-source link labeled ``original OWR report''. No individual problem page or passage is supplied. The citation uses 2017, whereas the DOI and proposed-year metadata refer to 2016; this distinction is retained without resolution.
\item Supplied literature-source link labeled ``transseries as surreal functions'': \url{https://arxiv.org/abs/1703.01995}. No authors, full bibliographic details, or specific supporting passages are supplied.
\item Literature assessment date: \texttt{2026-08-21}. This is supplied dated curation, not independently verified. The supplied status is ``open''. The supplied assessment states that the record collects several open problems and that composition on all of \(\mathbf{No}\), together with the listed elementary-equivalence and embedding questions, had no complete published resolution as of that date. The background describes this as literature triage derived from the supplied OWR curation and requests independent verification of time-sensitive status and cited claims.
\item The supplied statement-status label is \texttt{corrected\_verified}. Its accompanying verification note reports that the official report passage was checked and rewritten as a self-contained problem preserving its mathematical content. This is a supplied provenance claim, not verification performed for this draft; the missing definitions and discrepancy between the original and rewritten statements remain unresolved.
\end{itemize}
\end{document}
